% Source-aligned English for printed p.139.
\SourcePage{139}

\LayerTarget{tohoku:def:quotient-composition:p2}
\noindent the subobject $B''$ of $B$ to the quotient $B'$. The object $Q$ is
also isomorphic to the coimage of this morphism and is therefore both a
subobject of $B'$ and a quotient object of $B''$. By shrinking the subobject
$A'$ of $A$ and the quotient $C'$ of $C$ if necessary, we may suppose that
$u$ and $v$ come respectively from morphisms, still denoted by the same
letters, $A'\to Q$ and $Q\to C'$. We may therefore
take the composite $vu\in\mathrm{Hom}_{\mathrm{C}}(A',C')$. The element of
$\mathrm{Hom}_{\mathrm{C}/\mathrm{C}'}(A,C)$ that it defines depends only on $u'$ and $v'$; we
denote it by $v'u'$. It is straightforward to prove that the composition law
so defined is bilinear and associative, and that the class in
$\mathrm{Hom}_{\mathrm{C}/\mathrm{C}'}(A,A)$ of the identity map $I_A$ is an identity element.
Thus $\mathbf{C}/\mathbf{C}'$ is an additive category; in fact it is even an abelian category.
We shall not give this extremely tedious verification here.
\LayerTarget{tohoku:prop:quotient-category-properties}
Thus \textit{$\mathbf{C}/\mathbf{C}'$ is an abelian category; moreover the identity functor $F$
from $\mathbf{C}$ to $\mathbf{C}/\mathbf{C}'$ is exact} and in particular preserves kernels, cokernels,
images, and coimages. \textit{$F(A)$ is a zero object if and only if $A\in \mathbf{C}'$, and
every object of $\mathbf{C}/\mathbf{C}'$ is of the form $F(A)$ for some $A\in \mathbf{C}$.} These facts,
which essentially characterize the quotient category, permit the language
\textquotedblleft modulo $\mathbf{C}'$\textquotedblright\ to be used safely: it simply
means that one is working in the abelian quotient category.
\LayerTarget{tohoku:remark:quotient-spectral-sequences}
This language is particularly convenient when one has a spectral sequence in
$\mathbf{C}$ and wants to exploit the fact that some of its terms lie in $\mathbf{C}'$; see
2.4. Reducing modulo $\mathbf{C}'$, that is, applying the functor $F$, gives a spectral
sequence in $\mathbf{C}/\mathbf{C}'$ in which the corresponding terms vanish. One thereby
often obtains exact sequences modulo $\mathbf{C}'$ by the usual criteria for deriving
exact sequences from spectral sequences in which certain terms vanish.

\LayerTarget{tohoku:chapter:2}
\bigskip
\begin{center}
{\normalsize\bfseries Chapter II Homological Algebra in Abelian Categories}
\end{center}

\LayerTarget{tohoku:section:2.1}
\textbf{2.1. $\partial$-functors and
$\partial^*$-functors.}
\LayerTarget{tohoku:def:delta-functor:p1}
Let $\mathbf{C}$ be an abelian category and $\mathbf{C}'$ an additive category, and let $a$ and
$b$ be two integers, possibly equal to $+\infty$ or $-\infty$, such that
$a+1<b$. A \textit{covariant $\partial$-functor from $\mathbf{C}$ to $\mathbf{C}'$ in degrees
$a<i<b$} is a system $T=(T^i)$ of additive covariant functors from $\mathbf{C}$ to
$\mathbf{C}'$, one for each $a<i<b$, together with, for every $i$ such that
$a<i<b-1$ and every exact sequence
$0\to A'\to A\to A''\to 0$
in $\mathbf{C}$, a morphism

\[
 \partial:T^i(A'')\to T^{i+1}(A')
\]

\noindent called the boundary homomorphism, or connecting homomorphism,
satisfying the following axioms:

\LayerTarget{tohoku:axiom:delta-functor:1}
 (i) Given a second exact sequence
$0\to B'\to B\to B''\to 0$
and a homomorphism from the first exact sequence to the second, the
corresponding diagram

\LayerTarget{tohoku:diagram:delta-naturality}
\[
\begin{array}{lcl}
T^i(A'') & \mathop{\to}\limits^{\partial} & T^{i+1}(A')\\
\big\downarrow && \big\downarrow\\
T^i(B'') & \mathop{\to}\limits^{\partial} & T^{i+1}(B')
\end{array}
\]

\noindent is commutative.

\LayerTarget{tohoku:axiom:delta-functor:2:p1}
 (ii) For every exact sequence
$0\to A'\to A\to A''\to 0$,
the associated sequence of morphisms

\TohokuContinuation{140}
